% Generated draft: requires mathematical and source review.
\documentclass[11pt]{article}
\usepackage[margin=25mm]{geometry}
\usepackage{fontspec}
\setmainfont{Latin Modern Roman}
\usepackage{amsmath,amssymb,mathtools,unicode-math}
\setmathfont{Latin Modern Math}
\usepackage{enumitem,needspace,xurl,hyperref}
\setlength{\parindent}{0pt}
\setlength{\parskip}{5pt}
\setlength{\emergencystretch}{3em}
\newcommand{\N}{\mathbb{N}}
\newcommand{\Z}{\mathbb{Z}}
\newcommand{\Q}{\mathbb{Q}}
\newcommand{\R}{\mathbb{R}}
\newcommand{\C}{\mathbb{C}}
\newcommand{\F}{\mathbb{F}}
\newcommand{\A}{\mathbb{A}}
\newcommand{\PP}{\mathbb{P}}
\newcommand{\E}{\mathbb{E}}
\newcommand{\eps}{\varepsilon}
\newcommand{\dd}{\,\mathrm{d}}
\begin{document}
\section*{UnsolvedMath ID: 801166}
\textbf{Source problem code:} catalogue entry
\subsection{Definitions and mathematical statement}
Under the usual dyadic dilation and integer translation convention, let \(H=L^2(\R)\), with norm \(\|f\|_2=(\int_{\R}|f(x)|^2\,\dd x)^{1/2}\). A function \(\psi\in H\) is a Riesz wavelet if the family
\[
\psi_{j,k}(x)=2^{j/2}\psi(2^j x-k),\qquad j,k\in\Z,
\]
is a Riesz basis of \(H\): its closed linear span is \(H\), and there exist constants \(0<A\le B<\infty\) such that every finitely supported scalar family \((c_{j,k})\) satisfies
\[
A\sum_{j,k\in\Z}|c_{j,k}|^2
\le \left\|\sum_{j,k\in\Z}c_{j,k}\psi_{j,k}\right\|_2^2
\le B\sum_{j,k\in\Z}|c_{j,k}|^2.
\]
The constants may depend on \(\psi\). Let \(\mathcal W_R\) denote the collection of these Riesz wavelets.

The question attributed to Larson, reference label [La2], asks whether
\[
\overline{\mathcal W_R}^{\,\|\cdot\|_2}=L^2(\R).
\]
Equivalently, for every \(f\in L^2(\R)\) and every \(\varepsilon>0\), does there exist \(\psi\in\mathcal W_R\) such that \(\|f-\psi\|_2<\varepsilon\)?

The supplied excerpt does not define Riesz wavelets or specify the scalar field. The dyadic and integer translation convention above is an explicit interpretation requiring confirmation against the cited source; the scalar field of \(H\) and of the coefficients must agree with that intended there. The supplied registry labels the problem open; this label is not an independently verified literature assessment.
\subsection{Short English statement}
Can every function in \(L^2(\R)\) be approximated arbitrarily closely in norm by a single Riesz wavelet? The displayed definition uses the usual dyadic convention, which requires source confirmation. The source attributes the question to Larson.
\subsection{Sources}
\begin{itemize}
\item Local UnsolvedMath registry, permanent ID \(801166\), source problem code \texttt{catalogue entry}, titled \emph{Open Problems in Wavelet Theory --- Problem 5}. The supplied record lists Bownik and Rzeszotnik, proposed year 2019, and status \texttt{open}.
\item Bownik and Rzeszotnik, \emph{Open Problems in Wavelet Theory} (2019), Problem 5: \url{https://arxiv.org/abs/1905.08230}. The supplied provenance records access on 2026-10-16 and extraction by \texttt{arxiv\_latex}; these are supplied provenance claims, not independent source verification.
\item Larson, reference label [La2] in the supplied statement, is credited with posing the density question. No full bibliographic entry for [La2] is supplied.
\end{itemize}
\end{document}
